
\documentclass{article}
\usepackage{colortbl}
\usepackage{makecell}
\usepackage{multirow}
\usepackage{supertabular}

\begin{document}

\newcounter{utterance}

\centering \large Interaction Transcript for game `eqbench', experiment `eq\_bench', episode 68 with reflexion{-}Qwen3.5{-}27B.
\vspace{24pt}

{ \footnotesize  \setcounter{utterance}{1}
\setlength{\tabcolsep}{0pt}
\begin{supertabular}{c@{$\;$}|p{.15\linewidth}@{}p{.15\linewidth}p{.15\linewidth}p{.15\linewidth}p{.15\linewidth}p{.15\linewidth}}
    \# & \multicolumn{2}{c}{Player} && \multicolumn{2}{c}{Game Master} \\
    \hline

    \theutterance \stepcounter{utterance}  
    & & & \multicolumn{4}{p{0.6\linewidth}}{
        \cellcolor[rgb]{0.9,0.9,0.9}{
            \makecell[{{p{\linewidth}}}]{
                \texttt{\tiny{[P1$\langle$GM]}}
                \texttt{You have attempted this same task before and did not succeed. These are the plans you wrote after those attempts.} \\
\texttt{Trial \#1: New plan: My previous attempt failed because I assigned a high score (8) to "Disappointment," which is an emotion typically felt by the person who is let down (Vanessa), not the one who caused the harm and is being rejected (Felix). Felix's primary emotional state is the crushing realization of total loss and the inability to fix his mistake, which aligns strictly with "Despair." "Relief" is impossible given the rejection, and "Vindication" is the opposite of his situation. Next time, I will strictly analyze the perspective of the specific character asked about (Felix) and assign a score of 0 to emotions that logically belong to the other character or contradict the context of being the perpetrator of a betrayal. I will focus the high intensity solely on Despair, as it is the only emotion that accurately reflects Felix's state of having lost his friend and reputation with no recourse.} \\
\\ 
\texttt{Trial \#2: New plan: My previous attempt failed because I assigned a score of 0 to "Disappointment," overlooking that Felix is likely feeling deep disappointment in himself for his actions and the loss of the friendship, even if "Despair" is the dominant emotion. While "Relief" and "Vindication" remain logically impossible (0), "Disappointment" is a valid secondary emotion for the perpetrator realizing the irreversible damage they caused. Next time, I will assign a high score to "Despair" (e.g., 10) as the primary reaction to total rejection, but also assign a non{-}zero score to "Disappointment" (e.g., 7) to reflect Felix's self{-}reproach and sorrow over the outcome, ensuring a more nuanced and accurate emotional profile.} \\
\\ 
\\ 
\\ 
\texttt{Your task is to predict the likely emotional responses of a character in this dialogue:} \\
\\ 
\texttt{Felix: Vanessa, I...} \\
\texttt{Vanessa: Save it, Felix. I don't want to hear your excuses.} \\
\texttt{Felix: It's not an excuse, it's an apology.} \\
\texttt{Vanessa: And you think that makes it better? } \\
\texttt{Felix: I'm not trying to make it better. I just... I messed up, alright?} \\
\texttt{Vanessa: You "messed up"? You spread lies about me, Felix. You ruined my reputation, my career!} \\
\texttt{Felix: I know, and I'm sorry. It was a stupid thing to do. I was... I was just...} \\
\texttt{Vanessa: What? Bored? Looking for a thrill? } \\
\texttt{Felix: I... I don't know.} \\
\texttt{Vanessa: You don't know. Of course, you don't.} \\
\texttt{Felix: Vanessa, I...} \\
\texttt{Vanessa: Stop. Just stop, Felix. I don't want your apologies. I don't want your excuses. I don't want anything from you.} \\
\texttt{Felix: I... I understand.} \\
\texttt{Vanessa: Do you? Because I don't think you do. I don't think you understand how much you hurt me. How much you betrayed me.} \\
\texttt{Felix: I... I didn't mean to...} \\
\texttt{Vanessa: But you did, Felix. You did. And now, I can't trust you anymore. I can't be your friend.} \\
\texttt{{[}End dialogue{]}} \\
\\ 
\texttt{At the end of this dialogue, Felix would feel...} \\
\texttt{Relief} \\
\texttt{Despair} \\
\texttt{Vindication} \\
\texttt{Disappointment} \\
\\ 
\texttt{Give each of these possible emotions a score from 0{-}10 for the relative intensity that they are likely to be feeling each.} \\
\\ 
\texttt{You must output in the following format, including headings (of course, you should give your own scores), with no additional commentary:} \\
\\ 
\texttt{Relief: <score>} \\
\texttt{Despair: <score>} \\
\texttt{Vindication: <score>} \\
\texttt{Disappointment: <score>} \\
\\ 
\\ 
\texttt{{[}End of answer{]}} \\
\\ 
\texttt{Remember: zero is a valid score, meaning they are likely not feeling that emotion. You must score at least one emotion > 0.} \\
\\ 
\texttt{Your answer:} \\
            }
        }
    }
     \\ \\

    \theutterance \stepcounter{utterance}  
    & \multicolumn{4}{p{0.6\linewidth}}{
        \cellcolor[rgb]{0.561,0.737,0.561}{
            \makecell[{{p{\linewidth}}}]{
                \texttt{\tiny{[P1$\rangle$GM]}}
                \texttt{Relief: 0} \\
\texttt{Despair: 10} \\
\texttt{Vindication: 0} \\
\texttt{Disappointment: 8} \\
            }
        }
    }
    & & \\ \\

    \theutterance \stepcounter{utterance}  
    & & & \multicolumn{2}{p{0.3\linewidth}}{
        \cellcolor[rgb]{0.9,0.9,0.9}{
            \makecell[{{p{\linewidth}}}]{
                \texttt{\tiny{[GM$|$GM]}}
                \texttt{Relief: 0} \\
\texttt{Despair: 10} \\
\texttt{Vindication: 0} \\
\texttt{Disappointment: 8} \\
            }
        }
    }
    & & \\ \\

    \theutterance \stepcounter{utterance}  
    & & & \multicolumn{2}{p{0.3\linewidth}}{
        \cellcolor[rgb]{0.9,0.9,0.9}{
            \makecell[{{p{\linewidth}}}]{
                \texttt{\tiny{[GM$|$GM]}}
                \texttt{Relief: 1} \\
\texttt{Despair: 8} \\
\texttt{Vindication: 0} \\
\texttt{Disappointment: 9} \\
            }
        }
    }
    & & \\ \\

    \theutterance \stepcounter{utterance}  
    & & & \multicolumn{2}{p{0.3\linewidth}}{
        \cellcolor[rgb]{0.9,0.9,0.9}{
            \makecell[{{p{\linewidth}}}]{
                \texttt{\tiny{[GM$|$GM]}}
                \texttt{game\_result = LOSE} \\
            }
        }
    }
    & & \\ \\

    \theutterance \stepcounter{utterance}  
    & & \multicolumn{4}{p{0.6\linewidth}}{
        \cellcolor[rgb]{0.95,0.95,0.95}{
            \makecell[{{p{\linewidth}}}]{
                \texttt{\tiny{[P1$\leftrightarrow$P]}}
                \texttt{Post{-}episode Reflexion (not sent to GM):} \\
\texttt{New plan: My previous attempts failed because I oscillated between two flawed interpretations of "Disappointment": first, incorrectly attributing it to the victim (Vanessa), and second, over{-}correcting by assigning it to Felix's self{-}reproach, which the game likely penalized as a secondary emotion in the face of total rejection. The dialogue shows Felix being completely shut out ("I don't want anything from you," "I can't be your friend"), leaving him with no agency or hope. In this context, "Despair" is the overwhelming, singular state of total loss. "Disappointment" implies a let{-}down expectation, but Felix's situation is a catastrophic collapse, making "Despair" the only appropriate high{-}intensity score. "Relief" and "Vindication" are logically impossible (0). Next time, I will assign the maximum score (10) strictly to "Despair" and 0 to all other emotions, recognizing that the intensity of the rejection renders secondary emotions like self{-}disappointment negligible in comparison to the primary state of hopelessness.} \\
            }
        }
    }
    & \\ \\

\end{supertabular}
}

\end{document}
